\documentclass[11pt]{article}
\usepackage[a4paper, top=18mm, bottom=18mm, left=20mm, right=20mm]{geometry}
\usepackage[T2A]{fontenc}
\usepackage[utf8]{inputenc}
\usepackage[ukrainian]{babel}
\usepackage{graphicx}
\usepackage[export]{adjustbox}
\usepackage{amsmath}
\usepackage{amssymb}
\graphicspath{{/app/Quiz/Scripts/2023/ProgAll/15BinTree/}}
\setlength{\parindent}{0pt}
\setlength{\parskip}{6pt}
\emergencystretch=3em
\pagestyle{empty}
\begin{document}
%begin
{\centering\Large\bfseries Контрольна робота\par}
\medskip
\noindent\rule{\linewidth}{0.5pt}
\smallskip
\noindent\textbf{Варіант \No\,1}\hfill \textit{ПІБ:}\,\underline{\hspace{60mm}}
\vspace{4mm}

%beginitem
1. \textbf{Що буде виведено за виконання фрагменту коду?}

print('R', end=' ')
d=['0','1','2','3','4']
d.insert(1, [1,0])
print(d)

\par\vspace{5mm}
%enditem
%beginitem
2. 
Значенням змінної \kw{a} є цілочисловий масив типу $numpy.ndarray$ розмірів $6{\times}6$. На скільки збільшиться сума елементів масиву \kw{a} після виконання
\begin{verbatim}
a.T[0] += 3
\end{verbatim}


Якщо код некоректний, то в якості відповіді зазначити ERROR.


\par\vspace{5mm}
%enditem
%beginitem
3. \textbf{Що буде виведено за виконання фрагменту коду?}
%code
for d in range(-7,-7,2):
    if d>=6:
        break
        print(d,end=' ')
    if d>=6:
        break
else:
    print(d, end=' ')
print(d, end=' ')
%code
\par\vspace{5mm}
%enditem
%beginitem
4. %goodpath:Scripts/2019/Mod2019_1/Lex/01-good.txt
\textbf{Відмітити все, що є однією правильною лексемою:}
( 1 ) 1__23

( 2 ) 2j

( 3 ) True

( 4 ) 0123
\par\vspace{5mm}
%enditem
%beginitem
5. 

\begin{center}
	\adjustimage{max size={0.8\linewidth}{0.8\paperheight}}
	{../15BinTree/trees_exp/79.png}
\end{center}
Указати послідовність міток вершин дерева, зображеного на рис., що утвориться в результаті обходу дерева в оберненому порядку. Мітки записувати через пробіл.\par Алгоритм починає роботу з кореня (на рис. це верхня вершина).
%answer:79 оберненому
\par\vspace{5mm}
%enditem
%beginitem
6. \textbf{Що буде виведено за виконання фрагменту коду?}

a,b,c=3,4,8
def h(b):
global c    a*=4
    b=5
    c=2
    return a+b+c

a,b,c=6,7,1
print(h(b),a,b,c)
\par\vspace{5mm}
%enditem
%beginitem
7. \textbf{Що буде виведено за виконання фрагменту коду?}

%code
4**1.0-2
%code
\par\vspace{5mm}
%enditem
%beginitem
8. \textbf{Що буде виведено за виконання фрагменту коду?}

%code
cout << (15 % 15 % 15 % 15);
%code
\par\vspace{5mm}
%enditem
%beginitem
9. \textbf{Що буде виведено за виконання фрагменту коду?}

%code
print(5.0 >= 5 != False <= 6)
%code
\par\vspace{5mm}
%enditem
%beginitem
10. \textbf{Що буде виведено за виконання фрагменту коду?}
%code
a, b, c = 9, 6, 7
def h(a, b=8, c=6):
    print(a, b, c, end=" ")

h(3, 4, c=5)
print(a, b, c)
%code
\par\vspace{5mm}
%enditem










%end
\newpage

\newpage

%begin
{\centering\Large\bfseries Контрольна робота\par}
\medskip
\noindent\rule{\linewidth}{0.5pt}
\smallskip
\noindent\textbf{Варіант \No\,2}\hfill \textit{ПІБ:}\,\underline{\hspace{60mm}}
\vspace{4mm}

%beginitem
1. \textbf{Що буде виведено за виконання фрагменту коду?}

%code
#include <iostream>
using namespace std;

int h(int d){
    int x = 89;
    if (d >= 3) 
        x = 3;
    else if (d == 0)
         return 4;
    else
         return 8;
    return x;
}

int main(){
    cout << h(-1);
    return 0;
}
%code

\par\vspace{5mm}
%enditem
%beginitem
2. \textbf{Що буде виведено за виконання фрагменту коду?}

%code
print(4**1.0-2)
%code
\par\vspace{5mm}
%enditem
%beginitem
3. %goodpath:Scripts/2018/Mod2018_1/03-good.txt
\textbf{Відмітити все, що є коректним оператором С++:}
( 1 ) cin>>x+3;

( 2 ) (x<y && y>0)

( 3 ) return x<=2;

( 4 ) return *;
\par\vspace{5mm}
%enditem
%beginitem
4. \textbf{Що буде виведено за виконання фрагменту коду?}

%code
print('R', end=' ')
d = ('a','b','c','d','e',0,1,2,4)
print(d[-5])
%code

\par\vspace{5mm}
%enditem
%beginitem
5. \textbf{Що буде виведено за виконання фрагменту коду?}

%code
d = ('0','1','2','3','4','5','6','7','8','9')
print(d[-7:-7:2])
%code

\par\vspace{5mm}
%enditem
%beginitem
6. \textbf{Що буде виведено за виконання фрагменту коду?}
%code

a,b,c=9,6,7
def h(a,b=8,c=6):
    print(a,b,c,end="")

h(3,4,c=5)
print(a,b,c)

%code
\par\vspace{5mm}
%enditem
%beginitem
7. \textbf{Що буде виведено за виконання фрагменту коду?}

%code
cout << (!5.0>=5 || !false== 6 || 7!=6.0);
%code
\par\vspace{5mm}
%enditem
%beginitem
8. \textbf{Що буде виведено за виконання фрагменту коду?}

x,y,z=3,4,5
x,y,z,z=9,z,x,y
print(x,y,z)
\par\vspace{5mm}
%enditem
%beginitem
9. \textbf{Що буде виведено за виконання фрагменту коду?}
try:
    for d in range(-7, -7, 2):
        if d>=6:
            break
        print(d, end=';')
        if d>=6:
            break
    else:
        print(d,end=';')
    print(d)
except:
    print('ERR')
\par\vspace{5mm}
%enditem
%beginitem
10. \textbf{Що буде виведено за виконання фрагменту коду?}

%code
try:
    n = 13
    while n>=3:
        n += -2
    print(n, end=';')
except:
    print('Z')

%code
\par\vspace{5mm}
%enditem










%end
\newpage

\newpage

%begin
{\centering\Large\bfseries Контрольна робота\par}
\medskip
\noindent\rule{\linewidth}{0.5pt}
\smallskip
\noindent\textbf{Варіант \No\,3}\hfill \textit{ПІБ:}\,\underline{\hspace{60mm}}
\vspace{4mm}

%beginitem
1. \textbf{Що буде виведено за виконання фрагменту коду?}

%code
print('R', end=' ')
n=3
while n<=13:
    n+=2
print(n)
%code

\par\vspace{5mm}
%enditem
%beginitem
2. \textbf{Що буде виведено за виконання фрагменту коду?}
try:
    for d in range(-7, -7, 2):
        if d>=6:
            break
        print(d, end=';')
        if d>=6:
            break
    else:
        print(d,end=';')
    print(d)
except:
    print('Z')
\par\vspace{5mm}
%enditem
%beginitem
3. \textbf{Що буде виведено за виконання фрагменту коду?}
try:
    res = 6//6%6*6
    if isinstance(res, bool):
        print(f'bool;{res}')
    elif isinstance(res, int):
        print(f'int;{res}')
    elif isinstance(res, float):
        print(f'float;{res:.2f}')
    else:
        print('???')
except:
    print('Z')
\par\vspace{5mm}
%enditem
%beginitem
4. \textbf{Що буде виведено за виконання фрагменту коду?}

%code
cout << 15 % 15 % 15 % 15;
%code
\par\vspace{5mm}
%enditem
%beginitem
5. \textbf{Що буде виведено за виконання фрагменту коду?}
try:
    print(3, end=';')
    print(int('d4'), end=';')
    print(8, end=';')
except ZeroDivisionError:
    print(6, end=';')
except KeyboardInterrupt:
    print(7, end=';')
except:
    print('Z', end=';')
else:
    print(1, end=';')
finally:
    print(2)
\par\vspace{5mm}
%enditem
%beginitem
6. %goodpath:Scripts/2019/Mod2019_1/Lex/03-good.txt
\textbf{Відмітити всі коректні оператори Python:}
( 1 ) "a"=6

( 2 ) return x<=2

( 3 ) return 'a'<'c'

( 4 ) if a<b: a=0 else: b=0
\par\vspace{5mm}
%enditem
%beginitem
7. \textbf{Що буде виведено за виконання фрагменту коду?}
%code
#include <iostream>
using namespace std;

int f(int &x, int &y){
    x = 9;
    y-= 6;
    return y;
}

int main(){
    int a = 2, b = 3;
    b = f(b, b);
    cout << a << ":" << b <<':';
    {
        int a = 8, b = 5;
        cout << ((b>=6) || ((a-=2) >= 6))
             << a << ":" << b << ":";
    }
    {
        inta = 1;
        cout << a << ':';
    }
    cout << a << ":" << b << endl;
    return 0;
}
%code

\par\vspace{5mm}
%enditem
%beginitem
8. 
Вкажіть значення та тип виразу:
"9j">="9.0j"
\par\vspace{5mm}
%enditem
%beginitem
9. \textbf{Що буде виведено за виконання фрагменту коду?}
try:
    res = 4**1.0-2
    if isinstance(res, bool):
        print(f'bool;{res}')
    elif isinstance(res, int):
        print(f'int;{res}')
    elif isinstance(res, float):
        print(f'float;{res:.2f}')
    else:
        print('???')
except:
    print('Z')
\par\vspace{5mm}
%enditem
%beginitem
10. \textbf{Що буде виведено за виконання фрагменту коду?}
a = 3
b = 4
c = 8

def h(b):
global c    a *= 4
    b = 5
    c = 2
    return a + b + c

a = 6
b = 7
c = 1

try:
    print(h(b), a, b, c, sep=';')
except:
    print('Z', a, b, c, sep=';')
\par\vspace{5mm}
%enditem










%end
\newpage

\newpage

%begin
{\centering\Large\bfseries Контрольна робота\par}
\medskip
\noindent\rule{\linewidth}{0.5pt}
\smallskip
\noindent\textbf{Варіант \No\,4}\hfill \textit{ПІБ:}\,\underline{\hspace{60mm}}
\vspace{4mm}

%beginitem
1. 
Записати ВИРАЗ, значенням якого є список, що складається з кубів цілих чисел від 18 до 2008 (включно), що не діляться на жодне з чисел 2, 3, записаних в оберненому порядку.\par\vspace{5mm}
%enditem
%beginitem
2. \textbf{Що буде виведено за виконання фрагменту коду?}
try:
    print(3, end=';')
    print(int('d4'), end=';')
    print(8, end=';')
except ZeroDivisionError:
    print(6, end=';')
except KeyboardInterrupt:
    print(7, end=';')
except:
    print('ERR', end=';')
else:
    print(1, end=';')
finally:
    print(2)
\par\vspace{5mm}
%enditem
%beginitem
3. %goodpath:Scripts/2018/Mod2018_1/01-good.txt
\textbf{Відмітити все, що є однією правильною лексемою:}
( 1 ) +>

( 2 ) error

( 3 ) a2f

( 4 ) 0,5
\par\vspace{5mm}
%enditem
%beginitem
4. 
Написати програму, що запитує у користувача значення дійсних параметрів $x$ та $y$, обчислює для них значення виразу 
$$\frac{\sqrt x - 59}{(\ctg x - 0.9) (y - 66)}$$ 
та повiдомляє введенi данi та результат обчислення.

\par\vspace{5mm}
%enditem
%beginitem
5. \textbf{Що буде виведено за виконання фрагменту коду?}

print('R', end=' ')
for d in range(-7,-7,2):
    if d>=6:
        break
        print(d, end=' ')
    if d>=6:
        break
    else:
        print(d, end=' ')
    print(d)
\par\vspace{5mm}
%enditem
%beginitem
6. \textbf{Що буде виведено за виконання фрагменту коду?}

%code
6//6%6*6
%code
\par\vspace{5mm}
%enditem
%beginitem
7. 
Записати ВИРАЗ, значенням якого є множина, що складається з цілих чисел від 18 до 2008 (включно), що діляться на 7.\par\vspace{5mm}
%enditem
%beginitem
8. \textbf{Що буде виведено за виконання фрагменту коду?}
d=['0','1','2','3','4']; d.insert(1,[1,0]); print(d)\par\vspace{5mm}
%enditem
%beginitem
9. \textbf{Що буде виведено за виконання фрагменту коду?}
%Оригинал был утерян. Требует отладки!
%code
__d = 0

class D:
    __d = 1
    
    def __init__(self):
        self.__d = 2
        __d = 3
        D.__d = 4

obj = D()
D.__d = 5
try:
    print(__d, end=' ')
    print(obj.__d, end=' ')
    print(D.__d, end=' ')
    print(obj.__d, end=' ')
    print(obj._D__d)
except:
    print('ERR')
%code
\par\vspace{5mm}
%enditem
%beginitem
10. 
Написати програму, що запитує у користувача значення дійсних параметрів $x$ та $y$, обчислює для них значення виразу 
$$\frac{\sqrt x - 59}{(\ctg x - 0.9) (y - 66)}$$ 
та повiдомляє введенi данi та результат обчислення.
 Оцінюється правильність та якість коду.

\par\vspace{5mm}
%enditem










%end
\newpage
\end{document}


